\section*{Hoofdstuk \ref{ch:b2:mammal-reproduction} --- Geslachtelijke voortplanting van zoogdieren}

\begin{solution}{exo:b2:mammal-reproduction:1}
Spermatogonia delen zich vanaf de puberteit het hele leven lang; oögonia
alleen bij de foetus. Eén \omterm{def:b2:meiosis-heredity:meiosis}{meiose} geeft vier zaadcellen maar één eicel (en
drie poollichaampjes). Een zaadcel vergt ongeveer 64 dagen plus twee weken
rijping; een eicel staat stil in \omterm{prop:b2:meiosis-heredity:stages}{profase I} van vóór de geboorte tot de
maand van haar \omterm{def:b2:mammal-reproduction:ovary}{ovulatie} --- decennia --- en opnieuw in metafase II tot de
bevruchting. Een man maakt $10^{8}$ zaadcellen per dag, zo'n $10^{12}$ in
een leven; een vrouw ovuleert ongeveer 400 eicellen.
\end{solution}

\begin{solution}{exo:b2:mammal-reproduction:2}
Binding aan de \omterm{prop:b2:mammal-reproduction:fertilisation}{zona pellucida}; acrosoomreactie en het verteren van een
doorgang; versmelting van de membranen en binnenkomst van de kern van de
zaadcel; calciumgolf en \omterm{prop:b2:mammal-reproduction:fertilisation}{corticale reactie} (verharding van de zona, verlies
van haar receptoren: de blokkade van polyspermie); voltooiing van \omterm{def:b2:meiosis-heredity:meiosis}{meiose} II
en uitstoting van het tweede poollichaampje; vorming van de twee pronuclei,
DNA-replicatie, eerste mitose.
\end{solution}

\begin{solution}{exo:b2:mammal-reproduction:3}
\omterm{def:b2:mammal-reproduction:implantation}{Trofoblast} (buitenste laag) maakt de \omterm{def:b2:mammal-reproduction:placenta}{placenta} en de vliezen; \omterm{def:b2:mammal-reproduction:implantation}{binnenste
celmassa} maakt het embryo. De \omterm{def:b2:mammal-reproduction:implantation}{innesteling} begint op dag 6 tot 7 na de
bevruchting, in een baarmoederslijmvlies dat is voorbereid door
progesteron uit het \omterm{def:b2:mammal-reproduction:ovary}{corpus luteum}.
\end{solution}

\begin{solution}{exo:b2:mammal-reproduction:4}
Zuurstof: diffusie. Glucose: gefaciliteerd transport. Maternale rode
bloedcellen: passeren niet (de bloedsoorten mengen nooit). IgG:
receptorgemedieerd transport. Ureum: diffusie, van foetus naar moeder.
Alcohol: diffundeert vrij. Bacteriën: passeren meestal niet (sommige, zoals
die van syfilis en listeria, wel).
\end{solution}

\begin{solution}{exo:b2:mammal-reproduction:5}
$m = 3, 1, 0.4, 0.1$: $P = 0.95, 0.63, 0.33, 0.095$. Onder de helft
wanneer $m < \ln 2 = 0.69$, d.w.z.\ onder $7\times 10^{7}$ zaadcellen. De
curve verzadigt: boven ongeveer $10^{8}$ voegen extra zaadcellen vrijwel
niets toe, eronder daalt de vruchtbaarheid bijna evenredig met het aantal.
\end{solution}

\begin{solution}{exo:b2:mammal-reproduction:6}
Vóór de puberteit: $7\times 10^{5}$ verloren in 13 jaar (\num{4750} dagen):
ongeveer 150 per dag. Daarna: \num{299000} verloren in 37 jaar (\num{13500}
dagen): 22 per dag; \omterm{def:b2:mammal-reproduction:ovary}{ovulatie} $400/\num{13500} = 0.03$ per dag. Atresie is
zevenhonderd keer zo talrijk als \omterm{def:b2:mammal-reproduction:ovary}{ovulatie}; de voorraad raakt niet op door
ovuleren maar door afsterven.
\end{solution}

\begin{solution}{exo:b2:mammal-reproduction:7}
Volwassen: $(30/27)^{2.7} = 1.33$, $S = 1.33/2.33 = 0.57$. Foetaal:
$(30/19)^{2.7} = 3.4$, $S = 3.4/4.4 = 0.77$. Bij de lage drukken van de
\omterm{def:b2:mammal-reproduction:placenta}{placenta} laadt het foetale bloed zo'n twintig procentpunten meer zuurstof
dan volwassen bloed zou doen, en geeft het die in de weefsels af bij
drukken waarbij volwassen hemoglobine al half leeg zou zijn.
\end{solution}

\begin{solution}{exo:b2:mammal-reproduction:8}
$2^{k} = 25$: $k = 4.6$ verdubbelingen, \qty{9.3}{d}: dag 16 tot 17 na de
bevruchting; dag 30 tot 31 na de laatste menstruatie --- ongeveer de dag
waarop de menstruatie uitblijft.
\end{solution}

\begin{solution}{exo:b2:mammal-reproduction:9}
$3\times 10^{-8}\times 6\times 10^{4}\times 20/10^{-3} =
\qty{36}{mL/min}$ tegen een behoefte van \qty{25}{mL/min}: de marge daalt
van acht naar anderhalf; de foetus raakt bij elke extra belasting
zuurstoftekort, groeit traag en moet misschien vroeg worden geboren.
\end{solution}

\begin{solution}{exo:b2:mammal-reproduction:10}
(a) Er vormen zich testes die beide hormonen afscheiden; AMH verwijdert de
gangen van Müller, maar testosteron kan niet werken: de gangen van Wolff
verdwijnen en de uitwendige geslachtsorganen zijn vrouwelijk --- een vrouw
met testes en zonder baarmoeder. (b) Testosteron werkt, AMH niet: een man
met bijbal, zaadleider en mannelijke geslachtsorganen, plus een baarmoeder
en eileiders. (c) Ovaria en afgeleiden van de gangen van Müller (geen
\omterm{def:b2:mammal-reproduction:testis}{testis}, geen AMH); de androgenen uit de bijnier vermannelijken de
uitwendige geslachtsorganen in wisselende mate. (d)
\emph{SRY} maakt testes, die de mannelijke weg opleggen: een man met twee
X-chromosomen, onvruchtbaar omdat de Y-genen van de \omterm{def:b2:mammal-reproduction:testis}{spermatogenese}
ontbreken.
\end{solution}

\begin{solution}{exo:b2:mammal-reproduction:11}
De meest voorkomende (\qty{68}{\%}) is de splitsing op dag 4 tot 8: nadat
de \omterm{def:b2:mammal-reproduction:implantation}{trofoblast} is gevormd (dus één \omterm{def:b2:mammal-reproduction:placenta}{placenta}), maar vóór het amnion (dus twee
vruchtzakken). Vóór dag 3 maakt elke helft haar eigen \omterm{def:b2:mammal-reproduction:implantation}{trofoblast}, vandaar
twee \omterm{def:b2:mammal-reproduction:placenta}{placenta}'s (\qty{30}{\%}); na dag 8 is het amnion al aangelegd en
wordt het gedeeld (\qty{2}{\%}). Het splitsen van de \omterm{def:b2:mammal-reproduction:implantation}{binnenste celmassa}
nadat de \omterm{def:b2:mammal-reproduction:implantation}{trofoblast} is vastgelegd, is kennelijk de gemakkelijkste
gebeurtenis.
\end{solution}

\begin{solution}{exo:b2:mammal-reproduction:12}
Een buideldier investeert weinig vóór de geboorte: een korte dracht, een
piepklein jong, en een lactatie die kan worden afgebroken door het jong in
de buidel op te geven als er droogte komt, zodat het risico per poging
voor de moeder laag is en ze meteen opnieuw kan beginnen. Een placentaal
zoogdier investeert zwaar in een lange dracht met een groot, goed
ontwikkeld jong dat bij de geboorte beter overleeft, ten koste van een
moeder die maandenlang gebonden en intussen kwetsbaar is. In productieve,
voorspelbare omgevingen wint de hogere overleving van placentale jongen;
het grillige klimaat van Australië beloont de mogelijkheid van het
buideldier om op te geven, en zijn isolement hield placentale zoogdieren
tot voor kort buiten.
\end{solution}

\begin{solution}{pb:b2:mammal-reproduction:1}
\textbf{1.} $m = 3\times 10^{8}\times 10^{-8} = 3$; $P = 1 - e^{-3} =
0.95$.
\textbf{2.} $1 - e^{-3}(1 + 3) = 0.80$: zonder de blokkade zouden vier
bevruchtingen op vijf polysperm zijn en het embryo triploïd.
\textbf{3.} $0.15/5\times 10^{-5} = \qty{3000}{s}$, vijftig minuten;
de samentrekkingen van baarmoeder en eileider voeren ze mee.
\textbf{4.} $300/3\times 10^{8} = 10^{-6}$: verloren in de zure vagina, in
het slijm van de baarmoederhals, aan fagocyten in de baarmoeder, en in de
verkeerde eileider.
\textbf{5.} Van drie dagen vóór de \omterm{def:b2:mammal-reproduction:ovary}{ovulatie} tot één dag erna: ongeveer
vier dagen.
\textbf{6.} $m = 0.5$, $P = 0.39$; verwacht aantal cycli $1/P = 2.5$.
\textbf{7.} De weinige mitochondriën van de zaadcel, in het middenstuk,
worden na de binnenkomst gemerkt en vernietigd; de eicel draagt er
honderdduizend.
\textbf{8.} $120/16 = 7.5$: zeven delingen, $2^{7} = 128$ cellen ---
ongeveer honderd.
\textbf{9.} $\log_2 25 = 4.6$ verdubbelingen, \qty{9.3}{d}: dag 16 tot 17;
dag 30 tot 31 na de menstruatie.
\textbf{10.} Het \omterm{def:b2:mammal-reproduction:ovary}{corpus luteum} gaat ongeveer twaalf dagen na de \omterm{def:b2:mammal-reproduction:ovary}{ovulatie} in
regressie, tenzij hCG het redt; zonder hCG daalt het progesteron, wordt
het slijmvlies afgestoten en gaat het embryo met de menstruatie verloren.
\textbf{11.} $\log_2 10^{5} = 16.6$ verdubbelingen; tegen week 10 maakt de
\omterm{def:b2:mammal-reproduction:placenta}{placenta} haar eigen progesteron en is het \omterm{def:b2:mammal-reproduction:ovary}{corpus luteum} niet meer nodig.
\textbf{12.} Eicellen zitten decennialang in \omterm{prop:b2:meiosis-heredity:stages}{profase I} en de cohesines die
de bivalenten bijeenhouden, takelen af, zodat non-disjunctie bij \omterm{def:b2:meiosis-heredity:meiosis}{meiose} I
steil stijgt met de leeftijd van de moeder; zaadcellen worden in twee
maanden vers gemaakt.
\textbf{13.} Na de \omterm{def:b2:mammal-reproduction:implantation}{trofoblast} (dag 5) en vóór het amnion (dag
8): één \omterm{def:b2:mammal-reproduction:placenta}{placenta}, twee vruchtzakken.
\textbf{14.} $3\times 10^{-8}\times 1.2\times 10^{5}\times 20/3.5\times
10^{-4} = \qty{206}{mL/min}$.
\textbf{15.} $7\times 3.5 = \qty{24.5}{mL/min}$: een marge van acht.
\textbf{16.} $500\times 0.16 = \qty{80}{mL/min}$; extractie
$24.5/80 = \qty{31}{\%}$.
\textbf{17.} $5\times 3.5\times 1440 = \qty{25}{g/d}$; in \qty{28}{L}
maternaal bloed; dagelijkse bloedstroom \qty{720}{L}: de foetus neemt
ongeveer \qty{4}{\%} van de glucose die passeert.
\textbf{18.} Foetaal hemoglobine is bij \qty{30}{mmHg} voor \qty{77}{\%}
verzadigd, de foetus heeft meer rode bloedcellen en een hoog hartminuutvolume,
en zijn weefsels zijn aangepast aan werken bij lage druk --- hij leeft in
feite op de top van de Everest.
\textbf{19.} Gasuitwisseling zonder werkende long, voeding en uitscheiding
zonder werkende darm of nier, en de hormonen die de zwangerschap in stand
houden; ze kan niet elk toxine tegenhouden --- alcohol, nicotine en
sommige virussen passeren.
\textbf{20.} Rekking van de baarmoederhals $\to$ hypothalamus $\to$ oxytocine $\to$
samentrekking $\to$ meer rekking, tot de bevalling de prikkel wegneemt.
Vóór het einde van de zwangerschap heeft de door progesteron gedomineerde
baarmoeder weinig oxytocinereceptoren en geen gap junctions, zodat
samentrekkingen zich niet verspreiden en de lus niet kan sluiten.
\textbf{21.} $750\times 2.9 = \qty{2.2}{MJ}$; kosten $2.2/0.8 =
\qty{2.7}{MJ}$.
\textbf{22.} Groei $30\times 6 = \qty{0.18}{MJ}$, onderhoud
$300\times 4 = \qty{1.2}{MJ}$: \qty{1.4}{MJ}, ongeveer twee derde van de
energie van de melk; de rest gaat naar activiteit, warmte en verliezen.
\textbf{23.} $270\times 2.7 = \qty{730}{MJ}$, tweeënhalf keer de
zwangerschap.
\textbf{24.} Zogen verhoogt de prolactine en onderdrukt de pulsatiele
afgifte van het hormoon dat de \omterm{def:b2:mammal-reproduction:ovary}{ovulatie} aanstuurt, zodat de \omterm{def:b2:mammal-reproduction:ovary}{ovulatie} pas
maanden later terugkeert; zonder anticonceptie spreidt dit de geboorten
over twee tot drie jaar.
\textbf{25.} $P = 0.95$; test positief op dag 17 na de bevruchting
(dag 31 van de cyclus); zuurstofmarge acht; een dag melk kost de moeder
\qty{2.7}{MJ}.
\end{solution}
